\documentclass[10pt,a4paper]{amsart}
\usepackage[margin=19mm]{geometry}
\usepackage{amsmath,amssymb,mathtools}
\usepackage[scr=rsfs,cal=euler]{mathalfa}
\usepackage{xfrac}
\usepackage[unicode,hidelinks,bookmarks=false]{hyperref}
\newcommand{\ie}{i.e.\ }
\newcommand{\cf}{cf.\ }
\newcommand{\resp}{respectively\ }
\newcommand{\Subsection}{\subsection}
\providecommand{\nobreakdash}{\nobreak\hbox{-}}
\setlength{\emergencystretch}{2em}
\multlinegap0pt
\usepackage{amssymb}
\usepackage{mathtools}
\usepackage[shortlabels]{enumitem}
\newtheorem{theorem}{Theorem}
\title{On the category of semi-graded modules}
\author{Armando Reyes}
\date{Source: arXiv:2605.26374v1 (CC BY 4.0)}
\begin{document}
\maketitle

\noindent\emph{Source and licence.} On the category of semi-graded modules. Version arXiv:2605.26374v1 (CC BY 4.0), distributed by arXiv under the Creative Commons Attribution 4.0 International licence, http://creativecommons.org/licenses/by/4.0/, as recorded in the arXiv licence metadata for this exact version and shown on the abstract page https://arxiv.org/abs/2605.26374v1. Full licence notice: \texttt{licenses/arxiv-2605.26374-CC-BY-4.0.txt}.\par\smallskip

\noindent\textit{Editorial scope.} Counted unit: the article's theorem thm.BaersCriterion (source0.tex lines 261--266). The article's complete original proof of it is delivered with it (source0.tex lines 267--294, one \texttt{proof} environment of 28 lines). No mathematics is written, repaired or re-derived: the statement and the proof are the article's own lines, in the article's own notation, and no retained line is substituted or re-flowed.  Retained uncounted as the source-local context the proof needs: lines 205--207.  Nothing was omitted from any retained span, so the package carries no omission record.  No retained line contains a citation, so the wrapper carries no bibliography and no bibliography entry.  No retained line contains a figure environment, an image-inclusion command or drawing code, so no image is delivered and the package declares no asset dependency.  Editorial formatting only, in addition: an AMS \texttt{amsart} preamble that restates the article's own declarations and macros used by the retained lines, namely the environments \texttt{theorem} on separate counters, together with the standard packages those lines need.  The retained environments print in this document as Theorem 1 in order of appearance, whereas the article prints the same items with the numbering of its own full document.  The complete native source of the article as distributed in the arXiv e-print for this exact version is bundled unchanged as \texttt{sources/201412/source0.tex}.
\begin{theorem}[Grothendieck category]\label{Grothendieckcategory}
The family of objects $\left\{R(n)\right\}_{n \in \mathbb{Z}}$ constitutes a set of generators for the category $\mathsf{SGR}-R$. Consequently, $\mathsf{SGR}-R$ is a Grothendieck category and possesses enough injective objects.
\end{theorem}

\begin{theorem}[Baer's criterion]\label{thm.BaersCriterion}
Let $E$ be an object in $\mathsf{SGR}-R$. Then $E$ is an injective object if and only if for every $\mathsf{SG}$ left ideal $J \subseteq R$ and every strict homogeneous morphism of degree zero $g: J \to E$, there exists a homogeneous element $x \in E_0$ such that 
\[
g(r) = r x, \quad {\rm for\ all } \ r \in J.
\]
\end{theorem}

\begin{proof}
Assume that $E$ is an injective object in $\mathsf{SGR}-R$. Let $J \subseteq R$ be a $\mathsf{SG}$ left ideal, which means the canonical inclusion map $\iota: J \to R$ is a monomorphism of degree zero in $\mathsf{SGR}-R$. Let $g: J \to E$ be a degree-zero morphism. Since $h$ is a morphism in $\mathsf{SGR}-R$, it is $R$-linear and strictly maps $R_0 \to E_0$. Consider $x := h(1_R) \in E_0$. Using the left $R$-linearity of $h$, for any $r \in J \subseteq R$ we have 
\[
g(r) = (h \circ \iota)(r) = h(r) = h(r 1_R) = r h(1_R) = r x.
\]

Conversely, assume that every degree-zero morphism from any $\mathsf{SG}$ left ideal \linebreak $J \to E$ is given by right multiplication by an element in $E_0$. We must show that $E$ is an injective object. 

Let $A \subseteq B$ be a submodule inclusion in $\mathsf{SGR}-R$ and let $f: A \to E$ be a morphism of degree zero. 

We define a poset $\mathcal{P}$ of pairs $(A', f')$, where $A'$ is a $\mathsf{SG}$ submodule of $B$ containing $A$ ($A \subseteq A' \subseteq B$), and $f': A' \to E$ is a degree-zero morphism extending $f$ (i.e., $f'|_{A} = f$). We define a partial order on $\mathcal{P}$ as 
\[
(A'_1, f'_1) \preceq (A'_2, f'_2) \quad {\rm if \ and \ only \ if} \quad A'_1 \subseteq A'_2 \text{ and } f'_2|_{A'_1} = f'_1.
\]
Then every totally ordered chain $\left\{(A'_\alpha, f'_\alpha)\right\}_{\alpha \in \Omega}$ in $\mathcal{P}$ has an upper bound given by $\left(\bigcup_\alpha A'_\alpha, \bigcup_\alpha f'_\alpha\right)$. Since $\mathsf{SGR}-R$ is a Grothendieck category by Theorem \ref{Grothendieckcategory}, the directed union of $\mathsf{SG}$ submodules remains a well-defined $\mathsf{SG}$ submodule. By Zorn's Lemma, $\mathcal{P}$ contains a maximal element, say $(M, \psi)$. 

We claim that $M = B$. Suppose that there exists a homogeneous element $b \in B_j \setminus M$. We consider the {\em conductor ideal of $b$ into} $M$ as
\[
J := \{ r \in R \mid r b \in M \}.
\]
Because $M$ and $b$ are homogeneous, $J$ is a $\mathsf{SG}$ left ideal of $R$. We define a degree-zero map $\theta: J \to E$ by $\theta(r) = \psi(r b)$. Since $J$ is a $\mathsf{SG}$ ideal, there exists some $x \in E_0$ such that $\theta(r) = r x$ for all $r \in J$. Thus, $\psi(r b) = r x$.

Now, consider the larger semi-graded submodule $M' = M + R b$. We extend $\psi$ to a map $\psi': M' \to E$ by defining it component wise on homogeneous elements. For $u \in M_k$ and $r \in R_{k-j} = (R(-j))_k$, define 
\[
\psi'(u + r b) = \psi(u) + r x.
\]
To verify this map is well-defined, if $u + r b = 0$, then $r b = -u \in M$, which means $r \in J$. Then $r x = \theta(r) = \psi(r b) = \psi(-u) = -\psi(u)$, which shows $\psi(u) + r x = 0$. Since $\psi'$ is of degree zero and $R$-linear, $(M', \psi') \in \mathcal{P}$. However, $(M, \psi) \prec (M', \psi')$ because $b \notin M$, which contradicts the maximality of $(M, \psi)$. In this way, $M = B$, meaning $\psi: B \to E$ is the global extension we desire. From this, $E$ is an injective object.
\end{proof}

\end{document}
